\subsection{The prescribed imaginary-quadratic twist}
\label{sec:actual-cm-twist}

The rank-two seed uses an algebraic character with a specified
polarization and prime-to-$\ell$ finite inertia.  We construct that
character over the original field, with every field of the prescribed
contract retained~\cite{cm-twist-source}.

\begin{theorem}\label{thm:actual-cm-twist}
Let $K$ be any imaginary quadratic field of discriminant $D$, let
$\ell$ be odd with $\ell\nmid D$, and fix the original coefficient
isomorphism $\iota:\overline{\mathbb Q}_\ell\simeq\mathbb C$ and
complex conjugation $c$.  There is a continuous character
$\kappa:G_K\to\overline{\mathbb Q}_\ell^{\,\times}$ satisfying the
entire original \texttt{IsTwist} contract.  More explicitly, it is the
geometric-Artin avatar of an algebraic Hecke character $Y$ with
infinity type $(-1,0)$ and conductor supported on $D$, and
\[
 \kappa(s)\kappa(csc^{-1})=\varepsilon(s)^{-1},\qquad
 \kappa(t^2)=-\varepsilon(t)^{-1}
       \quad(s\in G_K,\ t\in G_{\mathbb Q}\setminus G_K).
\]
It is unramified outside $D\ell$.  At every finite place away from
$\ell$, its inertia image is killed by the integer
\[
 m_K=\begin{cases}54,&D=-3,\\32,&D\ne-3.\end{cases}
\]
In particular $\ell\nmid m_K$.  Its values lie in a finite extension
of $\mathbb Q_\ell$.  No restriction on the class number of $K$ is
required.
\end{theorem}

\begin{proof}
Let $S$ be the rational primes dividing $D$ and let $\eta$ be the
quadratic idele-class character of $\mathbb Q$ associated with
$K/\mathbb Q$.  Its infinity component is the sign character.
Fix a complex embedding $\tau:K\hookrightarrow\mathbb C$.

\par\medskip\noindent\textbf{Finite unit characters with the required order.}
We first construct
\[
 \nu=\prod_{p\in S}\nu_p:
       \prod_{\mathfrak p\mid D}\mathcal O_{K,\mathfrak p}^{\times}
                    \longrightarrow\mu_{m_K}
\]
such that $\nu_p$ restricts to $\eta_p$ on $\mathbb Z_p^\times$
and $\nu(u)=\tau(u)$ for every global unit $u$.
At an odd ramified prime other than the exceptional case $K=\mathbb
Q(\sqrt{-3})$, take the quadratic character of the residue field
$\mathbb F_p^\times$.  Its restriction is $\eta_p$: principal
units are squares when $p$ is odd, and the ramified quadratic
character is the nontrivial quadratic residue character on units.
Its order divides two.

At a ramified prime two, use the finite unit group
\[
                    U_2=(\mathcal O_K/8\mathcal O_K)^\times.
\]
It has order $32$, since the ring has $64$ elements and residue
field $\mathbb F_2$.  The rational unit character $\eta_2$ factors
through $(\mathbb Z/8)^\times$, because $1+8\mathbb Z_2$ consists
of squares.  Its embedded rational unit subgroup is literally
$(\mathbb Z/8)^\times$, since $\mathcal O_K$ is free over
$\mathbb Z$.  A character of a subgroup of a finite abelian group
extends to the whole group: adjoining one element $a$, choose a
root of the prescribed value of its first power lying in the
old subgroup.  Induction proves the assertion.  The resulting
character has values among the roots of unity of order dividing
the order of the whole group.  Thus $\eta_2$ extends to a
$\mu_{32}$-valued character of $U_2$.

For $K=\mathbb Q(i)$, prescribe on the larger subgroup generated
by rational units and $i$ the value $\nu_2(i)=\tau(i)$ before
extending.  These prescriptions agree on the intersection
$\{1,-1\}$, because $\eta_2(-1)=-1$.  The same intersection
statement holds modulo $8$, as is seen in the basis $1,i$.
For $K=\mathbb Q(\sqrt{-3})$, write $\zeta^2+\zeta+1=0$ and
instead use $(\mathcal O_K/9)^\times$, of order $54$.
Prescribe $\eta_3$ on rational units and $\tau$ on $\mu_6$.
Their intersection modulo $9$ is again exactly $\{1,-1\}$,
and $\eta_3(-1)=-1$; hence they extend to the whole finite group.
These prescriptions prove the global-unit condition in the two
exceptional fields.  In every other imaginary quadratic field
the only units are $\pm1$, and global reciprocity gives
\[
       \nu(-1)=\prod_{p\mid D}\eta_p(-1)=-1=\tau(-1).
\]
This constructs $\nu$ with its claimed order.  Let $\mathfrak f$
be the product of the chosen local moduli; it is supported on $D$.

\par\medskip\noindent\textbf{Extension across the actual ideal class group.}
Let $I_{\mathfrak f}$ be the group of fractional ideals prime to
$\mathfrak f$, and let $P_{\mathfrak f}$ denote its subgroup of
all principal ideals.  Thus a generator $a$ of an ideal in
$P_{\mathfrak f}$ is a unit at every prime in $\mathfrak f$;
no congruence $a\equiv1$ is imposed here.  Set
\[
             \psi_0((a))=\frac{\tau(a)}{\nu(a)}.
\]
The global-unit identity makes this independent of the generator.
The quotient $I_{\mathfrak f}/P_{\mathfrak f}$ is the finite
ordinary ideal class group: every ideal class has a representative
prime to a prescribed finite set, by approximation.

Decompose that finite quotient into independent cyclic factors.
Choose one generator in each factor, represented by an ideal
$\mathfrak a_i$ of order $h_i$, and write
$\mathfrak a_i^{h_i}=(a_i)$.  Choose once and for all an
$h_i$-th root of $\psi_0((a_i))$.  Assign these roots to the
$\mathfrak a_i$ and use the unique expression consisting of a
principal ideal times $\prod_i\mathfrak a_i^{e_i}$,
$0\le e_i<h_i$, to extend $\psi_0$ to a homomorphism
\[
                  \psi:I_{\mathfrak f}\longrightarrow
                                      \overline{\mathbb Q}^{\times}.
\]
The relations $\mathfrak a_i^{h_i}=(a_i)$ prove multiplicativity.
All values lie in one number field, obtained by adjoining the
finitely many chosen roots to $\tau(K)(\mu_{m_K})$.
An $\ell$-factor in a class-group order may enlarge this coefficient
field, but introduces no new finite inertia character.

Here is an explicit idele character, to keep its conductor visible.
For an idele $x$, choose $a\in K^\times$ such that $ax$ is a unit
at every prime in $\mathfrak f$, and let $I(ax)$ be its fractional
ideal outside that set.  Under the chosen complex identification put
\[
             Y(x)=(\tau(a)x_\infty)^{-1}
                         \nu((ax)_{\mathfrak f})\,\psi(I(ax)).
\]
Replacing $a$ by $ba$, with $b$ a unit at $\mathfrak f$,
multiplies the three factors by $\tau(b)^{-1}$, $\nu(b)$ and
$\tau(b)/\nu(b)$, respectively.  Thus $Y$ is well-defined.
It is multiplicative, continuous, and trivial on $K^\times$.
Its infinity component is $z^{-1}$, its unit component at a
bad prime is precisely $\nu_p$, and it is unramified at every
other finite prime.  A good local uniformizer has value
$\psi(\mathfrak q)$.

For every positive rational $b$ prime to $D$, quadratic reciprocity,
or the restriction of $\eta$ to the rational principal ideles, gives
\[
                \nu(b)=\eta_D(b),\qquad
                \psi((b))=b\eta_D(b),                 \tag{*}
\]
where $\eta_D$ is the primitive quadratic Dirichlet character of
discriminant $D$.  In particular, for $q\nmid D$,
$\psi((q))=q$ if $q$ splits, and $\psi((q))=-q$ if it is inert.
These identities hold for every choice made in extending the class
character.

\par\medskip\noindent\textbf{The original Galois normalization and polarization.}
Apply the algebraic-Hecke-character correspondence to $Y$
\cite{cm-twist-cht}.  The source's Artin map takes a uniformizer
to geometric Frobenius.  With infinity exponents $a=(-1,0)$,
the avatar is explicitly
\[
 \kappa(\operatorname{Art}_K(x))
 =\iota^{-1}\!\left(Y(x)\prod_\sigma x_\infty^{-a_\sigma}\right)
                         \prod_{\sigma\mid\ell}x_\sigma^{a_\sigma}.
\]
This is exactly the original \texttt{IsLAdicOf} formula.
The correspondence gives continuity and a finite coefficient field;
at $\mathfrak q\nmid D\ell$ it gives
\[
       \iota\bigl(\kappa(\operatorname{Frob}_{\mathfrak q}^{\rm ar})\bigr)
                                  =\psi(\mathfrak q)^{-1}.
\]
At every prime away from $\ell$, local reciprocity identifies its
inertia values with those of $\nu_p$, or with one at a good prime.
This proves both the original unramified clause and its single
prime-to-$\ell$ inertia exponent.  The trivial valuation has trivial
inertia, so the original valuation-subring formulation is included.

Consider the two-dimensional induced representation
$R=\operatorname{Ind}_{G_K}^{G_{\mathbb Q}}\kappa$.
At a split good prime $q=\mathfrak q\bar{\mathfrak q}$,
\[
       \det R(\operatorname{Frob}_q^{\rm ar})
         =(\psi(\mathfrak q)\psi(\bar{\mathfrak q}))^{-1}=q^{-1}.
\]
At an inert good prime, the induced matrix exchanges its two cosets,
so its determinant is
$-\kappa((\operatorname{Frob}_q^{\rm ar})^2)
 =-\psi((q))^{-1}=q^{-1}$ by (*).
Chebotarev and continuity therefore give
$\det R=\varepsilon^{-1}$ on all of $G_{\mathbb Q}$
\cite{cm-twist-cft}.  For every $t\notin G_K$, the same coset
matrix has determinant $-\kappa(t^2)$.  This is the required
\texttt{ver} identity for every such $t$, not only Frobenius elements.
Finally apply it to $t=sc$, using $c^2=1$ and
$\varepsilon(c)=-1$.  It gives
\[
      \kappa(s)\kappa(csc^{-1})
       =\kappa((sc)^2)=-\varepsilon(sc)^{-1}
       =\varepsilon(s)^{-1}.
\]
Every field of the original \texttt{IsTwist} record is now supplied.
\end{proof}

Restrict this very character to the original CM field used in the
rank-two seed.  Its source and normalization are unchanged.  This
discharges that seed's previously supplied CM-twist input; it does
not supply global marked deformation support, two-sided rigidity
or a uniform integral Selmer bound.
